\chapter{Giới thiệu}
\label{Chapter1}

\section{Lý do chọn đề tài}

Trong bối cảnh toàn cầu hóa hiện nay, năng lực sử dụng tiếng Anh đã trở thành một trong những tiêu chuẩn bắt buộc cho mọi hoạt động học tập và phát triển nghề nghiệp. Nền tảng của sự thông thạo ngôn ngữ này nằm ở việc xây dựng và duy trì một vốn từ vựng phong phú, bởi đây là yếu tố then chốt chi phối trực tiếp đến hiệu quả của cả bốn kỹ năng nghe, nói, đọc và viết \cite{Nation2001}. Tuy nhiên, thực trạng hiện nay cho thấy người học thường đối mặt với thách thức lớn trong việc chuyển dịch thông tin từ trí nhớ ngắn hạn sang dài hạn, dẫn đến sự đứt gãy trong lộ trình học tập do hiện tượng quên lãng tự nhiên.

Mặc dù các cơ sở khoa học về nhận thức như kỹ thuật lặp lại có khoảng cách (Spaced Repetition) và thực hành truy hồi (Retrieval Practice) đã chứng minh được hiệu quả vượt trội trong việc củng cố dấu vết trí nhớ \cite{Cepeda2008, Karpicke2008}, nhưng thực tế triển khai trên các nền tảng học ngôn ngữ phổ biến vẫn bộc lộ một khoảng trống quan trọng hơn: vấn đề duy trì động lực học tập dài hạn. Người học không bỏ cuộc vì thiếu phương pháp — họ bỏ cuộc vì quá trình lặp lại từ vựng ngày qua ngày trở nên đơn điệu, nhàm chán và thiếu cơ chế phản hồi đủ để duy trì động lực học tập trong dài hạn.

Anki cung cấp thuật toán lặp lại ngắt quãng chính xác nhưng giao diện tối giản và thiếu phản hồi tích cực. Người học dễ rơi vào trạng thái học cơ học thụ động — lật thẻ, đánh giá, lật thẻ tiếp — mà không nhận thấy sự tiến bộ hay gắn kết với nội dung, dẫn đến giảm động lực sau vài tuần đầu. Duolingo ở chiều ngược lại, cố gắng bù đắp bằng hệ thống động lực ngoại tại (Extrinsic Motivation): chuỗi ngày liên tục (streak), bảng xếp hạng điểm số, phần thưởng ảo. Song theo Thuyết Tự Quyết (Self-Determination Theory), các cơ chế ngoại tại này không nuôi dưỡng được hai yếu tố cốt lõi của động lực bền vững là cảm giác năng lực (Competence) và tính tự chủ (Autonomy). Khi streak bị gián đoạn hoặc khi phần thưởng ngoại tại mất đi sức hấp dẫn, toàn bộ động lực học tập sụp đổ theo, để lại người học không có lý do nội tại nào để tiếp tục.

Nhằm giải quyết những tồn tại này, khóa luận đã thực hiện đề tài \textbf{"Phát triển hệ thống hỗ trợ học từ vựng tiếng Anh sử dụng ASP.NET Core và ReactJS kết hợp gamification"}. Đề tài được lựa chọn dựa trên ba trụ cột mục tiêu chiến lược: 
(1) Hiện thực hóa quy trình ghi nhớ dựa trên thuật toán SM-2 tiêu chuẩn, bổ sung hàm chấm điểm đa yếu tố tự động (\texttt{CalculateSm2Grade}) để thay thế bước tự đánh giá thủ công trong SM-2 gốc bằng tính toán dựa trên độ chính xác, tốc độ phản hồi và số gợi ý sử dụng; 
(2) Xây dựng cơ chế gắn kết học tập thông qua hệ thống thú ảo gắn liền với tiến trình học tập; 
(3) Hiện thực hóa giải pháp trên một nền tảng kỹ thuật vững chắc theo kiến trúc hành tây (Onion Architecture) với hệ sinh thái .NET 8 và ReactJS để đảm bảo tính mở rộng và khả năng bảo trì cao.

\section{Mục tiêu đề tài}
Mục tiêu của đề tài là nghiên cứu và xây dựng một nền tảng học từ vựng trực tuyến, nhằm giải quyết hai thách thức trong việc tiếp thu ngôn ngữ: hiệu suất ghi nhớ dài hạn và khả năng duy trì động lực học tập. Để đạt được điều này, hệ thống hiện thực hóa thuật toán lặp lại ngắt quãng SuperMemo-2 (SM-2), một mô hình toán học tiên tiến cho phép cá nhân hóa hoàn toàn lộ trình ôn tập dựa trên khả năng phản hồi thực tế của từng cá nhân. Việc ứng dụng thuật toán này giúp tối ưu hóa "thời điểm vàng" để truy hồi kiến thức, từ đó hỗ trợ người học chuyển dịch thông tin từ trí nhớ ngắn hạn sang trí nhớ dài hạn một cách hiệu quả nhất.

Bên cạnh giải pháp về mặt kỹ thuật nhận thức, đề tài còn tập trung nghiên cứu sự giao thoa giữa tâm lý học hành vi và phát triển phần mềm thông qua việc tích hợp các cơ chế trò chơi hóa (Gamification) chuyên sâu. Hệ thống không chỉ dừng lại ở các yếu tố bề nổi mà còn xây dựng một cấu trúc tương tác phức hợp, tiêu biểu là cơ chế nuôi dưỡng và tiến hóa thú ảo. Phương pháp tiếp cận này nhằm tăng tính tương tác trong quá trình học tập, qua đó đáp ứng các nhu cầu tâm lý về sự tự chủ (Autonomy) và năng lực (Competence) theo Thuyết Tự Quyết (Self-Determination Theory). Sự kết hợp giữa cơ sở khoa học nhận thức và cơ chế trò chơi hóa được kỳ vọng hỗ trợ người học duy trì thói quen ôn tập từ vựng hiệu quả hơn.

\section{Phạm vi đề tài}

Đề tài tập trung vào việc nghiên cứu và xây dựng một ứng dụng web (Web Application) có khả năng vận hành đa nền tảng, hướng tới đối tượng mục tiêu là người học tiếng Anh ở trình độ từ cơ bản đến trung cấp. Phạm vi của nghiên cứu được xác định cụ thể thông qua sự kết hợp giữa các tính năng nghiệp vụ sư phạm và các tiêu chuẩn kỹ thuật phần mềm hiện đại.

Về mặt chức năng, hệ thống ưu tiên thiết lập các công cụ quản lý bộ từ vựng cá nhân hóa, cho phép người học chủ động trong việc xây dựng kho học liệu riêng. Quy trình học tập được hiện thực hóa thông qua bốn dạng tương tác cốt lõi bao gồm: Flashcard, Fill-in-blank (điền vào chỗ trống), Multiple Choice (trắc nghiệm) và Listening (nghe hiểu). Các phương thức này được thiết kế để bao quát các khía cạnh của trí nhớ ngôn ngữ từ nhận diện đến truy hồi kiến thức. Bên cạnh đó, đề tài chú trọng triển khai hệ thống quản lý thú ảo như một cơ chế trò chơi hóa chuyên sâu, đóng vai trò then chốt trong việc duy trì động lực và sự gắn kết của người học với nền tảng.

Về phương diện kỹ thuật, đề tài tập trung vào việc thiết kế và hiện thực hóa hệ thống theo kiến trúc hành tây (Onion Architecture), hay còn gọi là Clean Architecture. Việc áp dụng mô hình kiến trúc này nhằm đảm bảo sự tách biệt giữa logic nghiệp vụ (Business Logic) và tầng hạ tầng kỹ thuật (Infrastructure), từ đó cải thiện khả năng kiểm thử, bảo trì và tính độc lập giữa các lớp.

Cần lưu ý rằng hệ thống bao gồm nhiều tính năng (thú ảo, PvP/PvE, nhiệm vụ, thành tích, quản trị) vì các tính năng này được thiết kế như một hệ thống Gamification tích hợp, trong đó đóng góp kỹ thuật cốt lõi là hiện thực hóa thuật toán SM-2 và cơ chế ghi nhớ đa dạng; các tính năng Gamification đóng vai trò là phương tiện duy trì động lực học tập, không phải là mục tiêu nghiên cứu độc lập. Ngoài phạm vi trên, đề tài không bao gồm việc giảng dạy ngữ pháp chuyên sâu, luyện tập phát âm qua công nghệ nhận diện giọng nói (Speech Recognition), hay xây dựng các tính năng giao tiếp hội thoại tự do. Tương tác thời gian thực trong hệ thống được giới hạn trong phạm vi các phiên đối kháng trắc nghiệm có cấu trúc (PvP/PvE), không mở rộng sang các hình thức giao tiếp ngôn ngữ ngữ cảnh mở. Nhờ đó, nguồn lực được tập trung vào quy trình học tập và quản lý từ vựng.

\section{Phương pháp tiếp cận}

Đề tài được triển khai dựa trên sự kết hợp giữa nghiên cứu lý thuyết nền tảng và quy trình xây dựng phần mềm thực nghiệm chuyên nghiệp. Giai đoạn đầu của nghiên cứu tập trung vào việc tổng hợp và phân tích các cơ sở khoa học về nhận thức và tâm lý học hành vi, tiêu biểu là lý thuyết đường cong quên lãng của Ebbinghaus, thuật toán tối ưu hóa lịch trình ôn tập SM-2 và Thuyết Tự Quyết (Self-Determination Theory). Đây là những nền tảng quan trọng để thiết lập các quy tắc nghiệp vụ và logic vận hành cho hệ thống, đảm bảo tính khoa học trong việc hỗ trợ người học ghi nhớ từ vựng hiệu quả.

Về quy trình phát triển, dự án áp dụng mô hình Agile/Scrum rút gọn, cho phép phân tách quá trình thực hiện thành các giai đoạn ngắn để kiểm thử và tinh chỉnh liên tục. Phương pháp này phù hợp với việc phát triển các tính năng trò chơi hóa (Gamification), cho phép điều chỉnh theo kết quả kiểm thử thực tế.

Trong khâu thiết kế kiến trúc, hệ thống được mô hình hóa bằng ngôn ngữ UML để đảm bảo tính minh bạch về cấu trúc và luồng dữ liệu. Tầng Backend được tổ chức theo kiến trúc hành tây (Onion Architecture), kết hợp với các mẫu thiết kế (Design Patterns) tiêu biểu như Repository, Unit of Work và DTO. Việc lựa chọn kiến trúc này nhằm tách biệt logic nghiệp vụ khỏi hạ tầng kỹ thuật và cải thiện khả năng kiểm thử (Testability).

Hiện thực hóa hệ thống được thực hiện trên một bộ khung công nghệ hiện đại và đồng bộ. Tầng Backend tận dụng sức mạnh của .NET 8 và ASP.NET Core Web API, kết hợp với Entity Framework Core và SQL Server để quản lý dữ liệu quan hệ bền vững. Tầng Frontend sử dụng ReactJS và TypeScript để xây dựng giao diện người dùng tương tác cao, với sự hỗ trợ từ Vite và Tailwind CSS nhằm cải thiện tốc độ phát triển và hiệu suất giao diện. Cuối cùng, công tác kiểm thử và đánh giá được thực hiện nhiều tầng: Unit Test các logic nghiệp vụ cốt lõi, kiểm thử chức năng trên các kịch bản thực tế và kiểm thử hiệu năng tải (Load Testing) bằng công cụ k6 với kịch bản 100 người dùng ảo truy cập đồng thời, từ đó đối chiếu và xác nhận tính đúng đắn của các giả thuyết nghiên cứu đã đặt ra.

\section{Bố cục của khóa luận}

Khóa luận được tổ chức thành bốn chương chính, mỗi chương đảm nhận một giai đoạn riêng biệt trong lộ trình từ đặt vấn đề đến hoàn thiện sản phẩm.

\textbf{Chương 1 — Giới thiệu} xác lập bối cảnh nghiên cứu, trình bày lý do lựa chọn đề tài, hệ thống mục tiêu, phạm vi và phương pháp triển khai. Đây là khung tham chiếu chức năng và kỹ thuật định hướng cho toàn bộ công trình.

\textbf{Chương 2 — Cơ sở lý thuyết} tổng hợp và phân tích hệ thống nền tảng khoa học. Nội dung trọng tâm bao gồm cơ chế ghi nhớ lâu dài qua đường cong Ebbinghaus, mô hình toán học của thuật toán lặp lại ngắt quãng SM-2, các nguyên lý thiết kế Gamification hiệu quả dựa trên Thuyết Tự Quyết (SDT), và khảo sát các giải pháp liên quan hiện có trên thị trường.

\textbf{Chương 3 — Hệ thống hỗ trợ học từ vựng tiếng Anh kết hợp Gamification} là trọng tâm kỹ thuật của khóa luận. Chương này trình bày toàn bộ quy trình từ phân tích yêu cầu, thiết kế kiến trúc Onion Architecture, mô hình hóa hệ thống qua sơ đồ UML và lược đồ cơ sở dữ liệu quan hệ, cho đến việc trình bày kết quả hiện thực hóa các phân hệ chức năng cốt lõi và đánh giá hiệu năng hệ thống dưới tải.

\textbf{Chương 4 — Kết luận và Hướng phát triển} tổng kết toàn bộ công trình, đối chiếu kết quả đạt được so với mục tiêu ban đầu, nhìn nhận khách quan các hạn chế còn tồn tại và đề xuất lộ trình nâng cấp hướng đến một giải pháp EdTech hoàn chỉnh trong tương lai.